%==============================================================
%  lecons/analyse/wallis.tex — Intégrales de Wallis
%==============================================================

\lecontitle{Intégrales de Wallis}{Analyse réelle — Intégration et équivalents asymptotiques}

%=== NOTATION LOCALE =========================================
\newcommand{\Wn}{W_n}
\newcommand{\intz}{\int_0^{\pi/2}}

%=============================================================
%  § 1 — DÉFINITION ET ÉNONCÉ
%=============================================================
\secnum{navyblue}{1}{Définition et résultats fondamentaux}

\begin{tcolorbox}[thmbox]
  \textbf{Définition.}\enspace\index{intégrales de Wallis}
  Pour tout entier $n \geq 0$, on définit l'intégrale de Wallis d'ordre $n$ par
  \[
    \Wn \;=\; \intz \sin^n(t)\,\mathrm{d}t \;=\; \intz \cos^n(t)\,\mathrm{d}t.
  \]

  \medskip
  \textbf{Théorème (formules closes et équivalent).}
  \begin{itemize}
    \item \textit{Valeurs explicites :}
    \[
      W_{2p} = \frac{(2p-1)!!}{(2p)!!}\cdot\frac{\pi}{2}, \qquad
      W_{2p+1} = \frac{(2p)!!}{(2p+1)!!}.
    \]
    \item \textit{Relation de récurrence :}
    \[
      \Wn = \frac{n-1}{n}\,W_{n-2}, \quad W_0 = \frac{\pi}{2},\quad W_1 = 1.
    \]
    \item \textit{Équivalent asymptotique :}
    \[
      \Wn \;\sim\; \sqrt{\frac{\pi}{2n}} \quad (n \to +\infty).
    \]
    \item \textit{Produit de Wallis :}
    \[
      \prod_{k=1}^{+\infty} \frac{{(2k)}^2}{(2k-1)(2k+1)} \;=\; \frac{\pi}{2}.
    \]
  \end{itemize}
\end{tcolorbox}

\begin{significationintuitive}
Les intégrales de Wallis mesurent la « hauteur moyenne » de $\sin^n$ sur
$[0,\pi/2]$. Plus $n$ est grand, plus $\sin^n$ s'écrase vers zéro sauf au
voisinage de $\pi/2$ : $W_n$ tend vers $0$, mais lentement. L'équivalent
$W_n \sim \sqrt{\pi/(2n)}$ quantifie précisément cette décroissance en
$1/\sqrt{n}$, qui est à la racine de la formule de Stirling\index{formule de Stirling}.
\end{significationintuitive}

\decsep{}

%=============================================================
%  § 2 — DÉMONSTRATION
%=============================================================
\secnum{navyblue}{2}{Démonstration}

\begin{ideepreuve}
La récurrence s'obtient par intégration par parties. L'équivalent découle
de l'encadrement $W_{n+1} \leq W_n$ et du produit $W_{2p}\cdot W_{2p+1}$
calculé via les formules closes.
\end{ideepreuve}

\vspace{10pt}
\rubriq{Preuve}

\begin{mdframed}[style=proofbar]
  \textit{Étape 1 — Égalité sin/cos et monotonie.}
  Le changement $t \mapsto \pi/2-t$ montre l'égalité des deux intégrales.
  Comme $0\leq\sin t\leq 1$ sur $[0,\pi/2]$, la suite $(W_n)$ est décroissante.

  \medskip
  \textit{Étape 2 — Récurrence.}
  Pour $n\geq 2$, on intègre par parties avec $u = \sin^{n-1}t$ et
  $v' = \sin t$, donc $v = -\cos t$. Le crochet
  $\bigl[-\sin^{n-1}t\,\cos t\bigr]_0^{\pi/2}$ est nul ($\cos$ s'annule en
  $\pi/2$, $\sin^{n-1}$ en $0$), d'où, en écrivant $\cos^2 t = 1-\sin^2 t$ :
  \[
    \Wn = (n-1)\intz\cos^2 t\,\sin^{n-2}t\,\mathrm{d}t
         = (n-1)(W_{n-2}-\Wn),
  \]
  d'où $n\,\Wn = (n-1)\,W_{n-2}$, soit $\Wn = \dfrac{n-1}{n}W_{n-2}$.

  \medskip
  \textit{Étape 3 — Formules closes.}
  En déroulant depuis $W_0 = \pi/2$ (resp. $W_1 = 1$) :
  \[
    W_{2p} = \frac{(2p-1)!!}{(2p)!!}\cdot\frac{\pi}{2}, \qquad
    W_{2p+1} = \frac{(2p)!!}{(2p+1)!!}.
  \]

  \medskip
  \textit{Étape 4 — Équivalent.}
  De $W_{n+1}\leq W_n\leq W_{n-1}$ et $W_{n+1}/W_{n-1}=n/(n+1)\to 1$,
  on tire $W_n/W_{n-1}\to 1$. Les formules closes donnent :
  \[
    W_{2p}\cdot W_{2p+1} = \frac{\pi}{2(2p+1)}.
  \]
  Comme $W_{2p+1}\sim W_{2p}$, on obtient $W_{2p}^2\sim\pi/(4p)$, soit
  $W_{2p}\sim\sqrt{\pi/(4p)}$. Par symétrie des indices :
  \[
    \Wn\;\sim\;\sqrt{\frac{\pi}{2n}}.
  \]

  \medskip
  \textit{Étape 5 — Produit de Wallis.}
  Le rapport des formules closes se réécrit comme un produit télescopique :
  \[
    \frac{W_{2p}}{W_{2p+1}}
    = \frac{\pi}{2}\prod_{k=1}^{p}\frac{(2k-1)(2k+1)}{(2k)^2}.
  \]
  L'encadrement de l'étape~4 entraîne $W_{2p}/W_{2p+1}\to 1$ ; en passant à
  l'inverse, le produit converge et l'on obtient
  \[
    \prod_{k=1}^{+\infty}\frac{{(2k)}^2}{(2k-1)(2k+1)}=\frac{\pi}{2}.
  \]
  \hfill$\blacksquare$
\end{mdframed}

\decsep{}

%=============================================================
%  § 3 — EXEMPLES, CONTRE-EXEMPLES, EXERCICES
%=============================================================
\secnum{exgreen}{3}{Exemples, contre-exemples et exercices}

\rubriq{Exemples}

\begin{mdframed}[style=exbar]
  \textbf{1. Premières valeurs.}
  \[
    W_0=\tfrac{\pi}{2},\quad W_1=1,\quad W_2=\tfrac{\pi}{4},\quad
    W_3=\tfrac{2}{3},\quad W_4=\tfrac{3\pi}{16},\quad W_5=\tfrac{8}{15}.
  \]

  \smallskip\noindent
  \textbf{2. Lien avec Stirling.}
  Une fois établi par ailleurs (comparaison série-intégrale sur $\ln n!$)
  que $n!\sim C\sqrt{n}\,{(n/e)}^n$ pour une constante $C>0$, l'équivalent
  $W_{2p}^2\sim\pi/(4p)$, exprimé via $(2p)!$ et ${(p!)}^2$, identifie la
  constante $C=\sqrt{2\pi}$~: c'est la formule de Stirling
  $n!\sim\sqrt{2\pi n}\,{(n/e)}^n$ (voir exercice~3).

  \smallskip\noindent
  \textbf{3. Calcul direct.}
  $\displaystyle\int_0^{\pi/2}\sin^6 t\,\mathrm{d}t = W_6 = \tfrac{5}{6}\cdot\tfrac{3}{4}\cdot\tfrac{1}{2}\cdot\tfrac{\pi}{2} = \tfrac{5\pi}{32}$.

  \smallskip\noindent
  \textbf{4. Fonction Bêta.}
  $W_n = \tfrac{1}{2}B\!\left(\tfrac{1}{2},\tfrac{n+1}{2}\right)$, extension aux $n > -1$ réels.
\end{mdframed}

\vspace{6pt}
\rubriq{Contre-exemples et pièges}

\begin{mdframed}[style=cexbar]
  \textbf{Décroissance lente.}\enspace%
  $n\,W_n^2\to\pi/2\neq 0$ : la convergence vers $0$ est en $1/\sqrt{n}$,
  bien plus lente qu'une décroissance géométrique.

  \smallskip\noindent
  \textbf{Non-multiplicativité.}\enspace%
  $\Wn\neq W_{n-1}\cdot W_1$ : la relation de récurrence n'est pas un
  simple produit.
\end{mdframed}

\vspace{6pt}
\exolabel

\begin{mdframed}[style=exobar]
  \begin{enumerate}
    \item Calculer $\displaystyle\int_0^\pi\sin^7(t)\,\mathrm{d}t$ et
          $\displaystyle\int_0^{\pi/4}\cos^4(t)\,\mathrm{d}t$.

    \item Montrer que
          $\displaystyle\sum_{n=0}^{+\infty}{(-1)}^n W_{2n+1}
          =\int_0^{\pi/2}\frac{\sin t}{1+\sin^2 t}\,\mathrm{d}t
          =\frac{1}{2\sqrt{2}}\ln\!\frac{\sqrt{2}+1}{\sqrt{2}-1}\approx 0{,}623$.
          \textit{Indication :} sommer la série géométrique
          $\sum {(-1)}^n\sin^{2n+1}t$, puis justifier l'interversion
          somme-intégrale (la série n'étant pas absolument sommable) en
          majorant le reste de la série alternée par son premier terme~:
          $\bigl|\sum_{k>N}{(-1)}^k\sin^{2k+1}t\bigr|\leq\sin^{2N+3}t$,
          d'intégrale $W_{2N+3}\to 0$.

    \item En utilisant $W_n\sim\sqrt{\pi/(2n)}$, retrouver la formule
          de Stirling (exprimer $W_{2n}$ en termes de $(2n)!$ et ${(n!)}^2$).

    \item Montrer que $\displaystyle\int_0^{+\infty}e^{-t^2}t^{2n}\,\mathrm{d}t
          = \dfrac{(2n-1)!!}{2^{n+1}}\sqrt{\pi}$.

    \item (Généralisation) Pour $\alpha>-1$, poser
          $I_n(\alpha)=\int_0^{\pi/2}\sin^\alpha(t)\cos^n(t)\,\mathrm{d}t$.
          Établir une récurrence et exprimer $I_n(\alpha)$ via la fonction
          Bêta d'Euler.
  \end{enumerate}
\end{mdframed}

\leconfooter{J.\ Wallis, \textit{Arithmetica Infinitorum} (1656)}
